\section{Verification purpose}

This case defines a parametric verification of the resolution-independent WUI response to generation rate, inter-structure separation, and wind speed. The two controlled response families contain 270 requested points with duplicate baselines removed.

The response test is valid only when the integrated-structure capability is available. The current source gates that capability, so this case verifies design completeness and reports the limitation explicitly instead of substituting the native cell-local model.

\section{Mathematical formulation and numerical implementation}

Each physical structure is a square of side \(a=\SI{10}{m}\). Although the
computational grid uses \(\Delta x=\SI{2.5}{m}\), all cells belonging to one
structure must act as one source and one collector. For physical structure
\(s\), the required formulation is
\[
 \mathrm{HRR}_s(t)=q^{\prime\prime}(t)a^2,
 \qquad \dot N_s(t)=GR\times\mathrm{HRR}_s(t),
 \qquad N_s(t)=\sum_{j\in s}N_j(t).
\]
Transport is evaluated once for the structure-scale heat-release rate (HRR), deposition is
accumulated over its physical footprint, and the resulting load
\(N_s/a^2\) drives the physical ignition criterion. This grouping is the key
difference from the native multiple-discrete-cell treatment, in which one
physical structure is fragmented into independently emitting and igniting
raster cells.

The level-set time step is
\[
 \Delta t=0.5\frac{\Delta x}{U_{\rm wind}},
\]
and surface spread is disabled. Spread success therefore means that the
firebrand process alone propagates through the structural array. For each
parameter point, mean rate of spread (ROS) is obtained by fitting structure position
against first-ignition time.

The first family varies
\[
 d=2.5,5,\ldots,30\ \mathrm{m},
 \qquad GR=5,10,\ldots,50\
 \mathrm{pcs\,MW^{-1}\,s^{-1}},
\]
with a fixed 40~mph wind. The second fixes the separation at 10~m, uses the
same generation rates, and varies wind speed from 10 to 150~mph in 10~mph
increments.

\section{Derivation of expected behavior}
Increasing separation moves targets into lower-probability transport tails, while increasing generation raises deposited load. Wind produces competing reach and ignition-cooling effects, so the critical generation boundary should rise at both low and high wind speeds.

% BEGIN EXPLICIT EXPECTED-RESULT DERIVATION
The design count follows directly from the generated parameter products:
\[
 12\ \text{separations}\times10\ GR=120,\qquad
 15\ \text{winds}\times10\ GR=150,\qquad
 N_{\mathrm{design}}=120+150=270.
\]
The two baseline-equivalent coordinates remain distinct members of their
respective response families.  At the design-fire plateau, a physical
structure has \(HRR_{\max}=400(100)=40~\mathrm{MW}\); consequently the tested
generation range corresponds to
\[
 \dot N_{\max}=GR(40)=200,\ 400,\ldots,2000~\mathrm{pcs\,s^{-1}}.
\]
With \(\Delta x=2.5~\mathrm{m}\) and \(CFL=0.5\), the generated level-set
step is \(0.5(2.5)/[U_{\mathrm{mph}}(0.44704)]\): it is
0.069904 s at 40 mph, 0.279617 s at 10 mph, and 0.018641 s at 150 mph.

For each coordinate, \(GR_c\) is selected from
\(\{5,10,\ldots,50\}\); if no point propagates, the rank-only sentinel is the
next design level, 55.  Thus a one-level wind-edge excess is
\(55-50=5~\mathrm{pcs\,MW^{-1}\,s^{-1}}\), exactly the acceptance
increment.  The moderate set is 30--120 mph and the edge set is
\(\{10,20,130,140,150\}\) mph.  The explicit wind metrics are
\[
 \overline{GR}_{c,\mathrm{edge}}
 -\operatorname{median}(GR_{c,\mathrm{moderate}})\geq5,\qquad
 \frac{\overline{ROS}_{\mathrm{moderate}}}
      {\overline{ROS}_{\mathrm{edge}}}\geq1.10.
\]
The separation expectation is strong positive rank association of \(d\) with
\(GR_c\), negative rank association of \(d\) with ROS, and positive rank
association of GR with ROS.  These become
\(\rho\geq0.85\), median \(\rho\leq-0.50\), and median
\(\rho\geq0.70\), respectively, with the independent response-scale ceiling
\(\max ROS\leq0.08~\mathrm{m\,s^{-1}}\).
% END EXPLICIT EXPECTED-RESULT DERIVATION

\section{Verification metrics and acceptance criteria}

For a coordinate value \(z\), define the critical generation rate
\[
 GR_{c}(z)=\min\{GR:\text{the array sustains downwind spread}\}.
\]
When no tested value succeeds, the postprocessor assigns the next design
level, \(55\ \mathrm{pcs\,MW^{-1}\,s^{-1}}\), only for rank and boundary
metrics; it does not reinterpret failure as an observed threshold.

For the separation family, the documented response predicts a nonlinear
increase of \(GR_{c}\) with \(d\), from approximately 2.5 to more than
\(30\ \mathrm{pcs\,MW^{-1}\,s^{-1}}\)~\cite{Qin2025,QinEtAl2026}. ROS
should decrease with separation
and increase with generation. The verification therefore requires Spearman
rank correlation \(\rho(d,GR_{c})\ge0.85\), median
\(\rho(d,ROS)\le-0.50\), and median \(\rho(GR,ROS)\ge0.70\).

For the wind family, the critical boundary is U-shaped. Low wind gives
insufficient reach and a high ignition threshold; very high wind improves
transport but suppresses the smoldering-to-flaming transition through
convective cooling. The mean critical generation rate outside 30--120 mph
must exceed the median within that interval by at least
\(5\ \mathrm{pcs\,MW^{-1}\,s^{-1}}\). Mean ROS in the moderate range must
be at least 1.10 times the edge-wind mean. Finally, firebrand-driven ROS must
remain below \SI{0.08}{m.s^{-1}}, consistent with the documented response
scale~\cite{Qin2025,QinEtAl2026}.

\subsection{Justification of metric quantification}

The rank limits quantify robust trends without assuming that the nonlinear
responses are linear.  A separation-threshold Spearman coefficient of 0.85
requires strong ordering across twelve separation levels while allowing a
small number of adjacent ties or reversals near a discrete ignition boundary.
The \(-0.50\) ROS/separation and 0.70 ROS/generation limits require,
respectively, at least a moderate negative and a strong positive monotone
association despite flat-response regions.  They therefore test the direction
and persistence of the mechanisms rather than a fitted slope chosen for one
parameter interval.

The \(5\ \mathrm{pcs\,MW^{-1}\,s^{-1}}\) wind-edge excess equals one full
generation-rate design increment, so the U-shaped boundary must be resolved by
the experiment rather than inferred from sub-grid interpolation.  The 1.10
moderate/edge ROS ratio requires a 10\% effect beyond equality.  The
\SI{0.08}{m.s^{-1}} ceiling is a scale check taken from the reported
firebrand-driven response range for this design~\cite{Qin2025,QinEtAl2026}.
All 270 points and the resolution-independent structure capability are
mandatory because missing edge or threshold points can materially change both
rank statistics and the inferred U shape.

\subsection{Predeclared acceptance criteria}

Table~\ref{tab:acceptance-contract} summarizes the metrics fixed before
inspection of the calculated results. Numerical values and component statuses
are reported separately in Section~7.

\begin{table}[H]
\centering
\footnotesize
\caption{Predeclared verification metrics and acceptance criteria.}
\label{tab:acceptance-contract}
\begin{tabular}{@{}p{0.23\linewidth}p{0.47\linewidth}p{0.21\linewidth}@{}}
\toprule
Metric & Output, region, and aggregation & Acceptance criterion\\
\midrule
Design completeness & Unique generation--separation and generation--wind points after baseline deduplication. & 270 of 270 points \\
Required capability & Resolution-independent single-structure source, collector, and input data conversion procedure. & Available \\
Result completeness & Current admissible result rows divided by 270. & 100\% \\
Threshold/separation rank & Spearman \(\rho(d,GR_{c})\). & \(\ge0.85\) \\
ROS/separation rank & Median Spearman \(\rho(d,ROS)\) across generation levels. & \(\le-0.50\) \\
ROS/generation rank & Median Spearman \(\rho(GR,ROS)\) across both response families. & \(\ge0.70\) \\
Wind-edge threshold excess & Mean critical generation outside 30--120 mph minus the median inside that range. & \(\ge\SI{5}{pcs.MW^{-1}.s^{-1}}\) \\
Moderate/edge ROS ratio & Mean ROS at 30--120 mph divided by mean ROS outside that range. & \(\ge1.10\) \\
Maximum ROS & Maximum firebrand-driven ROS over all admissible design points. & \(\le\SI{0.08}{m.s^{-1}}\) \\
Overall decision & Capability, complete results, and every trend/scale metric are conjunctive. & PASS only if all pass \\
\bottomrule
\end{tabular}
\end{table}

\section{Simulation configuration and predicted outcome}

Although execution remains capability-gated, every design point now carries an
auditable spatial configuration. Structures are rasterized as a periodic
two-dimensional lattice over a \(220\times100\)-m physical domain. The
generation--separation family applies its selected gap in both coordinate
directions; the generation--wind family uses the \SI{10}{m} baseline gap.
Every footprint in the first upwind structure column is marked by the
non-positive initial level-set field.

Figure~\ref{fig:input-configuration} shows the baseline-separation geometry.
The two-cell numerical halo is excluded. The left panel verifies repetition in
both directions, while the right panel verifies that ignition spans the entire
first structure column instead of selecting only the lower-left structure.

\begin{figure}[H]
\centering
\EvidenceFigure[width=0.98\textwidth,height=0.42\textheight,keepaspectratio]{../figures/input_configuration.pdf}
\caption{Representative prepared two-dimensional structure lattice and first-column
initial ignition mask for the capability-gated parametric design.}
\label{fig:input-configuration}
\end{figure}

The Courant--Friedrichs--Lewy (CFL) number measures the distance
advected during one timestep relative to the grid spacing and
constrains the time discretization.

\begin{table}[H]
\centering\small
\caption{Controlled parametric design.}
\begin{tabular}{lll}
\toprule
Parameter & Separation family & Wind family\tabularnewline
\midrule
Grid spacing & \SI{2.5}{m} & \SI{2.5}{m}\tabularnewline
Physical domain & \(220\times100\) m & \(220\times100\) m\tabularnewline
Structure side & \SI{10}{m} & \SI{10}{m}\tabularnewline
Separation & 2.5--\SI{30}{m}, step \SI{2.5}{m} & \SI{10}{m}\tabularnewline
Wind & \SI{40}{mph} & 10--\SI{150}{mph}, step \SI{10}{mph}\tabularnewline
Generation & 5--50, step 5 & 5--50, step 5\tabularnewline
Spatial treatment & single physical unit & single physical unit\tabularnewline
Initial ignition & complete first structure column & complete first structure column\tabularnewline
CFL based on wind & 0.5 & 0.5\tabularnewline
Surface spread & disabled & disabled\tabularnewline
Number of points & 120 & 150\tabularnewline
\bottomrule
\end{tabular}
\end{table}

Expected ignition classifications form an increasing boundary in the
\((d,GR)\) plane and a U-shaped boundary in the
\((U_{\rm wind},GR)\) plane. The ROS maps should be largest for high
production, short separation, and moderate wind. Hamada-model curves included
in the documented parameter comparison~\cite{Qin2025,QinEtAl2026} are
contextual because Hamada includes radiation and flame contact; they are not
acceptance references for this firebrand-only test.

\subsection{Preprocessing, execution, and postprocessing}

The parameter design contains all 270 prescribed combinations. Each defines
co-registered structure-fraction, structure-identity, fuel-model, and initial
level-set fields. A combination is simulated only if the tested ELMFIRE version
supports the resolution-independent structure formulation and its required
input representation. Substituting the native cell-based formulation would
test a different discretization and would not establish the intended result.

Analysis requires the specified quantities for the complete parameter design,
derives critical boundaries, calculates the rank and response metrics, and
constructs the classification and spread-rate comparisons. Missing evidence
remains NOT EVALUABLE; an unavailable model capability is not replaced by a
different physical or numerical formulation.

\section{Actual simulation results}

No whole-domain ELMFIRE result is available because the required integrated-structure
capability is absent and no simulation condition is runnable under the current source interface.
The postprocessor therefore does not manufacture a raster or substitute an analytical
field for model output.  The case remains NOT EVALUABLE, and the response map below is
reported only as the prescribed verification envelope.

\CompletedVariantCount{} of \RequestedVariantCount{} exact-design result
points are present. The decision is \textbf{\OverallStatus}. Figure~\ref{fig:parametric}
shows only the expected qualitative boundaries while the capability is absent.
Every panel is labeled accordingly; no guide curve is represented as an
ELMFIRE result.

\begin{figure}[H]
\centering
\EvidenceFigure[width=\linewidth]{parametric_response_maps.pdf}
\caption{Parametric-response verification layout. The upper panels reproduce the
coordinate systems of the two parametric reference maps; the lower panels are
reserved for the corresponding ROS curves. In the current capability-gated
state, shaded regions and dashed lines communicate expected behavior only.}
\label{fig:parametric}
\end{figure}

\section{Calculated metrics and verification decision}

\begin{table}[H]
\centering\small
\caption{Calculated verification metrics and component decisions.}
\begin{tabular}{@{}p{0.37\linewidth}p{0.18\linewidth}p{0.17\linewidth}p{0.14\linewidth}@{}}
\toprule
Metric & Acceptance limit & Calculated & Status\tabularnewline
\midrule
\MetricRows
\bottomrule
\end{tabular}
\end{table}

The current decision is \textbf{\OverallStatus}, not a failure of the
firebrand equations. The intended resolution-independent parameter sweep
cannot be executed because the current source does not implement the
single-structure
source and collector required by the cited model description~\cite{Qin2025}. The native model would
split each 10-m structure into sixteen 2.5-m cells and change HRR, transport,
and accumulation; those outputs are not admissible evidence for this case.

The complete parameter design, expected behavior, result configuration definition, metrics,
figure generator, and standalone report are now present. Once the source
feature and its input data conversion procedure are implemented, the capability gate should be
enabled, all 270 points executed, and the same postprocessor used without
changing acceptance limits.

\section{Technical interpretation}
The capability gate prevents cell-local structural discretization from being presented as evidence for a single-structure parametric model. Once admissible results exist, boundary and rank metrics distinguish physical response trends from isolated successful points.

\section{Scope and limitations}
The design covers idealized, uniform arrays and verifies qualitative model response over the declared parameter grid. It remains not evaluable until the required single-structure source and collector are implemented.

\begin{thebibliography}{9}
\bibitem{Qin2025}
Y. Qin, \emph{A Physics-Based Eulerian Framework for Modeling Firebrand
Showering in Regional-Scale Wildland and Wildland--Urban Interface Fire
Simulations}, Ph.D. dissertation, University of Maryland, College Park, 2025.

\bibitem{QinEtAl2026}
Y. Qin et al., ``Simulations of firebrand-driven fire spread in
landscape-scale Wildland--Urban-Interface and urban conflagration models,''
\emph{Fire Safety Journal}, vol.~162, 104686, 2026,
doi:10.1016/j.firesaf.2026.104686.
\end{thebibliography}
